%%==============================================================================
%%  國立臺灣大學碩博士論文 LaTeX 範本  —  National Taiwan University Thesis
%%  ── 從這裡開始 (start here) ──
%%
%%  Compile with XeLaTeX (required for Chinese):
%%      latexmk -xelatex main.tex
%%  or:  xelatex main ; bibtex main ; xelatex main ; xelatex main
%%
%%  Class options:  master | doctor  ,  oneside | twoside
%%==============================================================================
\documentclass[master,oneside]{ntuthesis}

% Bibliography style — author–year is common in the sciences; swap for any
% journal .bst if your field prefers numeric citations.
\usepackage[round,authoryear]{natbib}
\bibliographystyle{plainnat}

% Add your own packages here (amsmath is preloaded below as a demo):
\usepackage{amsmath}

%------------------------------------------------------------------------------
%  ★ 全部論文資訊只改這裡一個地方 (edit ALL metadata in this one block)
%------------------------------------------------------------------------------
\ntusetup{
  degree        = master,                 % master | doctor
  collegeCH     = 理學院,
  collegeEN     = College of Science,
  deptCH        = 地質科學系,               % ← 改成你的系所 (change to your dept)
  deptEN        = Department of Geosciences,
  titleCH       = 國立臺灣大學碩博士論文 LaTeX 範本,
  titleEN       = A LaTeX Thesis Template for National Taiwan University,
  authorCH      = 王小明,
  authorEN      = Xiao-Ming Wang,
  studentID     = R12XXX001,
  advisorCH     = 指導教授姓名,
  advisorEN     = Advisor Name,
  advisortitleCH= 博士,
  advisortitleEN= Ph.D.,
  % coadvisorCH = 共同指導教授姓名,          % 有共同指導教授才填
  % coadvisorEN = Co-advisor Name,
  dateCH        = 中華民國一一五年七月,
  dateEN        = July 2026,
  examdateCH    = 一一五年七月一日,
  examdateEN    = {July 1, 2026},
  keywordsCH    = {臺灣大學, 學位論文, LaTeX, 範本, 排版},          % 規範要求 5–7 個
  keywordsEN    = {National Taiwan University, thesis, LaTeX, template, typesetting},
}

\begin{document}

%============================ FRONT MATTER ====================================
\maketitlepage        % 封面／書名頁 (附件1–2)
\makespine            % 書背／側邊 (§2) — 亦可用 bookspine.tex 單獨編譯後移除此行
\makecertificate      % 口試委員會審定書 (附件3) — 正式版以掃描簽名頁替換

\frontmatter          % 前置部分改用小寫羅馬頁碼
% 誌謝 Acknowledgements — regulation §4: optional, keep to about one page.
\begin{acknowledgements}
在此感謝指導教授的悉心指導，以及師長、實驗室夥伴與家人在論文完成過程中的
協助與支持。（此頁為謝辭範例，請自行改寫；原則上以不超過一頁為宜。）
\end{acknowledgements}
   % 謝辭 (§4，可選)
% 中文摘要 — regulation §5: ≤ 3 頁，含論述重點、方法或程序、結果與討論及結論。
% 關鍵詞（5–7 個）由 main.tex 的 \ntusetup{keywordsCH=...} 設定。
\begin{abstractCH}
本範本示範如何以 XeLaTeX 依《國立臺灣大學碩、博士學位論文格式規範》排版學位論文。
使用者只需於 \texttt{main.tex} 的 \texttt{\textbackslash ntusetup} 區塊填入系所、題目、
姓名、指導教授、關鍵詞與日期，即可自動產生符合規範的封面、書背、口試審定書、
中英文摘要、目次／圖次／表次與內文版面。（此為摘要範例，請替換為你的實際內容；
摘要不超過三頁，內容應涵蓋論述重點、方法或程序、結果與討論及結論。）
\end{abstractCH}
        % 中文摘要 (§5)
% English abstract — regulation §6: no more than three pages.
% Keywords (5–7) are set via \ntusetup{keywordsEN=...} in main.tex.
\begin{abstractEN}
This template demonstrates how to typeset an NTU master's thesis or doctoral
dissertation with XeLaTeX according to the university's official format
regulation. Fill in your department, title, name, advisor, keywords, and dates
in the \texttt{\textbackslash ntusetup} block of \texttt{main.tex}, and the
compliant cover, book spine, acceptance certificate, bilingual abstracts,
tables of contents/figures/tables, and body layout are generated automatically.
(Replace this placeholder with your actual abstract; it should not exceed three
pages and should cover the main argument, methods, results, and conclusions.)
\end{abstractEN}
        % 英文摘要 (§6)
% % 符號說明 — optional list of symbols/abbreviations. Enable in main.tex.
\begin{denotation}[3cm]
  \item[NTU] 國立臺灣大學 (National Taiwan University)
  \item[$n$] 樣本數 (sample size)
  \item[$\alpha$] 顯著水準 (significance level)
\end{denotation}
       % 符號說明 (可選)
\maketoc              % 目次／圖次／表次

%============================== MAIN MATTER ===================================
\mainmatter           % 正文改用阿拉伯頁碼
\chapter{緒論 Introduction}
\label{ch:intro}

本章示範中英文混排、章節標題、交叉參照與文獻引用。中文採 12 號楷書、1.5 倍行距，
英文採 12 號 Times New Roman，皆由 \texttt{ntuthesis.cls} 依規範自動設定。

\section{研究背景 Background}
段落中可自由混用中文與 English text。跨語言自動斷行已由 \texttt{xeCJK} 處理。
引用文獻使用 \texttt{natbib}：括號引用 \citep{knuth1984texbook}、
文中引用 \citet{lamport1994latex}。

\section{文獻引用與交叉參照 Citations and cross-references}
交叉參照範例：本論文結構見第\ref{ch:conclusion}章；方法見第\ref{ch:method}章。
公式、圖、表亦可交叉參照，見式\eqref{eq:demo}、圖\ref{fig:demo}、表\ref{tab:demo}。

\section{論文架構 Thesis organization}
\begin{itemize}
  \item 第\ref{ch:intro}章：緒論。
  \item 第\ref{ch:method}章：研究方法（示範圖、表、公式）。
  \item 第\ref{ch:conclusion}章：結論。
\end{itemize}

\chapter{研究方法 Methods}
\label{ch:method}

本章示範學位論文常用的三種元素：數學公式、圖、表。

\section{數學公式 Equations}
行內公式如 $E = mc^2$；編號的展示公式如式\eqref{eq:demo}：
\begin{equation}
  \label{eq:demo}
  \int_{a}^{b} f(x)\,\mathrm{d}x = F(b) - F(a).
\end{equation}

\section{圖 Figures}
以 \texttt{\textbackslash includegraphics} 置入圖檔（存於 \texttt{figures/}）。此處以
一個佔位方框代替真實圖檔，讓範本不需附圖即可編譯；實務上請改用註解中的寫法。
\begin{figure}[htbp]
  \centering
  % 真實用法：\includegraphics[width=0.7\textwidth]{figures/your-figure.pdf}
  \fbox{\begin{minipage}[c][4cm][c]{0.7\textwidth}\centering（範例圖檔佔位 placeholder figure）\end{minipage}}
  \caption{圖題範例。若圖擷取自參考文獻，須於圖題標註來源（規範八）。}
  \label{fig:demo}
\end{figure}

\section{表 Tables}
表格建議使用 \texttt{booktabs} 的三線表樣式，見表\ref{tab:demo}。
\begin{table}[htbp]
  \centering
  \caption{表題範例。若表擷取自參考文獻，須於表中標註來源（規範九）。}
  \label{tab:demo}
  \begin{tabular}{lcc}
    \toprule
    項目 Item & 數值 Value & 單位 Unit \\
    \midrule
    樣本數 $n$      & 120  & --   \\
    平均值 $\bar{x}$ & 3.14 & cm   \\
    標準差 $s$      & 0.27 & cm   \\
    \bottomrule
  \end{tabular}
\end{table}

\chapter{結論 Conclusions}
\label{ch:conclusion}

條列研究的主要發現與後續建議。將此範本內容替換為你的實際論文即可。

本範本已依《國立臺灣大學碩、博士學位論文格式規範》設定版面（邊界、字體、行距、
頁碼、封面、書背、審定書、摘要、目次），可作為臺大各系所碩博士生撰寫學位論文的
起點。


%============================== REFERENCES ====================================
\clearpage
\addcontentsline{toc}{chapter}{參考文獻}
\bibliography{back/references}

%============================== APPENDIX ======================================
\appendix
\chapter{附錄 Appendix}
\label{app:demo}

附錄可放置完整資料表、補充推導、程式碼、問卷等。附錄章節會以英文字母編號
（附錄 A、附錄 B …）。


\end{document}
